% Part:sets-functions-relations
% Chapter: size-of-sets
% Section: non-enumerability

\documentclass[../../../include/open-logic-section]{subfiles}

\begin{document}

\olfileid[tr]{sfr}{siz}{nen}

\olsection{\printtoken{S}{nonenumerable} Kümeler}

\begin{editorial}
  Bu kesim, \olref[enm]{sec}'teki tanımı kullanarak $\Bin^\omega$ ve
  $\Pow{\PosInt}$'ın sayılamazlığını ispatlar. \olref[enm-alt]{sec}'teki
  sürümden biraz daha temel ve ayrıntılı olacak biçimde tasarlanmıştır.
\end{editorial}

Pozitif tam sayılar kümesi $\PosInt$ gibi bazı kümeler sonsuzdur. Şimdiye
kadar gördüğümüz sonsuz küme örneklerinin hepsi !!{enumerable} idi. Ancak bu
özelliği taşımayan sonsuz kümeler de vardır. Böyle kümelere
\emph{!!{nonenumerable}} denir.

Öncelikle !!{nonenumerable} kümelerin varlığı belki şimdiden şaşırtıcıdır.
Her !!{enumerable} $A$ kümesi için bir !!{surjective} $f\colon\PosInt\to A$
fonksiyonu vardır. Bir küme !!{nonenumerable} ise böyle bir fonksiyon yoktur.
Yani $\PosInt$'ın sonsuz sayıdaki !!{element}ını $A$'ya eşleyen hiçbir
fonksiyon $A$'nın tamamını tüketemez. Dolayısıyla $A$'nın, sonsuz sayıdaki
pozitif tam sayıdan ``daha çok'' !!{element}ı vardır.

Bir kümenin !!{nonenumerable} olduğu nasıl ispatlanır? Böyle örten bir
fonksiyonun var olamayacağını göstermeniz gerekir. Buna eşdeğer olarak
$A$'nın elemanlarının tek yönlü sonsuz bir listede numaralandırılamayacağını
göstermelisiniz. Bunu yapmanın en iyi yolu, $A$'nın !!{element}larından oluşan
her listenin en az bir elemanı dışarıda bıraktığını ya da hiçbir
$f\colon\PosInt\to A$ fonksiyonunun !!{surjective} olamayacağını
göstermektir. Bunu Cantor'un \emph{köşegen yöntemi}yle yapabiliriz. $A$'nın
!!{element}larından oluşan $x_1$, $x_2$, \dots{} biçiminde bir liste verildiğinde,
kuruluşu gereği bu listede bulunamayacak başka bir $A$ elemanı kurarız.

İlk örneğimiz, boşluksuz bütün sonsuz $0$ ve $1$ dizilerinin kümesi
$\Bin^\omega$'dır.

\begin{thm}
\ollabel{thm:nonenum-bin-omega}
$\Bin^\omega$ !!{nonenumerable}dır.
\end{thm}

\begin{proof}
Tersini varsayarak $\Bin^\omega$'nın !!{enumerable} olduğunu, yani
$\Bin^\omega$'nın bütün !!{element}larından oluşan $s_1$, $s_2$, $s_3$,
$s_4$, \dots{} listesinin bulunduğunu varsayalım. Bu $s_i$'lerin her biri
başlı başına sonsuz bir $0$ ve $1$ dizisidir. Bu listedeki $i$'inci dizinin
$j$'inci elemanına $s_i(j)$ diyelim. O hâlde $i$'inci dizi $s_i$ şöyledir:
\[
s_i(1), s_i(2), s_i(3), \dots
\]

Bu listeyi ve içindeki her $s_i$ dizisinin elemanlarını bir dizi içinde
düzenleyebiliriz:
\[
\begin{array}{c|c|c|c|c|c}
& 1 & 2 & 3 & 4 & \dots \\\hline
1 & \mathbf{s_{1}(1)} & s_{1}(2) & s_{1}(3) & s_1(4) & \dots \\\hline
2 & s_{2}(1)& \mathbf{s_{2}(2)} & s_2(3) & s_2(4) & \dots \\\hline
3 & s_{3}(1)& s_{3}(2) & \mathbf{s_3(3)} & s_3(4) & \dots \\\hline
4 & s_{4}(1)& s_{4}(2) & s_4(3) & \mathbf{s_4(4)} & \dots \\\hline
\vdots & \vdots & \vdots & \vdots & \vdots & \mathbf{\ddots}
\end{array}
\]
Yandaki etiketler $s_1$, $s_2$, \dots{} listesindeki dizilerin sıra numaralarını,
üstteki sayılar ise tek tek dizilerin !!{element}larını numaralandırır.
Örneğin $s_1(1)$, $s_1$ dizisinin ilk !!{element}ı olan sayının, yani bir
$0$ ya da bir~$1$'in adıdır; diğerleri de böyledir.

Şimdi bu listede bulunamayacak sonsuz bir $0$ ve $1$ dizisi
$\overline{s}$'yi kuruyoruz. $\overline{s}$'nin tanımı $s_1$, $s_2$,
\dots{} listesine bağlı olacaktır. Sonsuz $0$ ve $1$ dizilerinden oluşan her
sonsuz liste, listede yer almadığı güvence altında olan sonsuz bir
$\overline{s}$ dizisi doğurur.

$\overline{s}$'yi tanımlamak için bütün !!{element}larının ne olduğunu, yani
bütün $n\in\PosInt$ için $\overline{s}(n)$'yi belirtiriz. Bunu yukarıdaki
dizinin köşegeni boyunca aşağı doğru okuyarak (``köşegen yöntemi'' adı buradan
gelir), sonra her $1$'i $0$ ve her $0$'ı~$1$ yaparak gerçekleştiririz. Daha
soyut bir ifadeyle, köşegenin $n$'inci !!{element}ı $s_n(n)$'nin $1$ ya da
$0$ olmasına göre $\overline{s}(n)$'yi $0$ ya da $1$ olarak tanımlarız:
\[
\overline{s}(n) =
\begin{cases}
1 & \text{$s_{n}(n) = 0$ ise}\\
0 & \text{$s_{n}(n) = 1$ ise.}
\end{cases}
\]
Durumlara göre tanımlar yerine formülleri tercih ediyorsanız
$\overline{s}(n)=1-s_n(n)$ olarak da tanımlayabilirsiniz.

$\overline{s}$, tablomuzun köşegenindeki bitlerin ters çevrilmesiyle elde
edilen sonsuz bir $0$ ve $1$ dizisidir. Dolayısıyla $\overline{s}$,
$\Bin^\omega$'nın bir !!a{element}ıdır. Fakat $s_1$,
$s_2$, \dots{} listesinde bulunamaz. Neden?

Listedeki ilk dizi $s_1$ olamaz; çünkü ilk !!{element}da $s_1$'den farklıdır.
$s_1(1)$ ne olursa olsun $\overline{s}(1)$'i onun tersi olarak tanımladık.
Listedeki ikinci dizi olamaz; çünkü $\overline{s}$ ikinci elemanda $s_2$'den
farklıdır: $s_2(2)$ sıfırsa $\overline{s}(2)$ birdir, birse sıfırdır. Bu
böyle sürer.

Daha kesin olarak, $\overline{s}$ listede olsaydı $\overline{s}=s_k$ olacak
biçimde bir $k$ bulunurdu. İki dizi, ancak ve ancak her konumda uyuşuyorsa
özdeştir; yani her $n$ için $\overline{s}(n)=s_k(n)$ olmalıdır. Dolayısıyla
özellikle $n=k$ alındığında $\overline{s}(k)=s_k(k)$ olması gerekirdi.
$s_k(k)$ ya $0$ ya da~$1$'dir. Sıfırsa $\overline{s}(k)$ bir olmalıdır;
$\overline{s}$'yi böyle tanımladık. Fakat $s_k(k)=1$ ise yine
$\overline{s}$'yi tanımlama biçimimiz gereği $\overline{s}(k)=0$'dır. İki
durumda da $\overline{s}(k)\neq s_k(k)$ olur.

$\Bin^\omega$'nın !!{element}larından oluşan bir $s_1$, $s_2$, \dots{}
listesinin bulunduğunu varsayarak başladık. Bu listeden, listede
bulunamayacağını ispatladığımız bir $\overline{s}$ dizisi kurduk. Fakat bütün
$s_i$'ler $0$ ve $1$ dizileriyse $\overline{s}$ de kesinlikle bir $0$ ve
$1$ dizisidir; yani $\overline{s}\in\Bin^\omega$'dır. Bu özellikle,
$\Bin^\omega$'nın \emph{bütün} !!{element}larından oluşan hiçbir listenin
olamayacağını gösterir: böyle herhangi bir listeden, listede yer almadığı
güvence altında olan bir $\overline{s}$ dizisi daha kurabiliriz. Dolayısıyla
$\Bin^\omega$'daki bütün dizilerin bir listesi olduğu varsayımı çelişkiye
götürür.
\end{proof}

\begin{explain}
Bu ispat yöntemine ``köşegenleştirme'' denir; çünkü $\overline{s}$'yi
tanımlamak için dizinin köşegenini kullanır. Köşegenleştirme bir dizinin
varlığını gerektirmez: gerçek bir dizi ve köşegen söz konusu olmadığında da
benzer bir fikirle kümelerin !!{enumerable} olmadığını gösterebiliriz.
\end{explain}

\begin{thm}
\ollabel{thm:nonenum-pownat}
$\Pow{\PosInt}$ !!{enumerable} değildir.
\end{thm}

\begin{proof}
Aynı biçimde ilerleyerek $\PosInt$'ın alt kümelerinden oluşan her liste için listede
bulunamayacak bir $\PosInt$ alt kümesinin olduğunu gösteririz. Aşağıdakinin
$\PosInt$ alt kümelerinin verilmiş bir listesi olduğunu varsayın:
\[
Z_{1}, Z_{2}, Z_{3}, \dots
\]
Şimdi her $n\in\PosInt$ için $n\in\overline{Z}$ ancak ve ancak $n\notin Z_n$
olacak biçimde bir $\overline{Z}$ kümesi tanımlarız:
\[
\overline{Z} = \Setabs{n \in \PosInt}{n \notin Z_n}.
\]
Varsayım gereği her $Z_n$ bir pozitif tam sayılar kümesi olduğundan
$\overline{Z}$ de açıkça bir pozitif tam sayılar kümesidir; dolayısıyla
$\overline{Z}\in\Pow{\PosInt}$. Fakat $\overline{Z}$ listede bulunamaz.
Bunu göstermek için her $k\in\PosInt$ için $\overline{Z}\neq Z_k$ olduğunu
göstereceğiz.

Öyleyse $k\in\PosInt$ gelişigüzel seçilsin. $\overline{Z}$'yi her
$n\in\PosInt$ için $n\in\overline{Z}$ ancak ve ancak $n\notin Z_n$ olacak
biçimde tanımladık. Özellikle $n=k$ alındığında $k\in\overline{Z}$ ancak ve
ancak $k\notin Z_k$ olur. Fakat bu, $k$ birinin !!a{element}ı olup diğerinin
olmadığından ve dolayısıyla $\overline{Z}$ ile $Z_k$ farklı !!{element}lara
sahip olduğundan $\overline{Z}\neq Z_k$ olduğunu gösterir. $k$ gelişigüzel
olduğuna göre $\overline{Z}$, $Z_1$, $Z_2$, \dots{} listesinde değildir.
\end{proof}

\begin{explain}
Önceki ispat köşegenden söz etmedi; ama şöyle canlandırırsanız bir köşegen
içerdiğini düşünebilirsiniz: $Z_1$, $Z_2$, \dots{} kümelerinin bir dizide
yazıldığını ve her $j\in Z_i$ için bu !!{element}ı $j$'inci sütuna yazdığımızı
düşünün. Listedeki ilk dört kümenin $\{1,2,3,\dots\}$, $\{2,4,6,\dots\}$,
$\{1,2,5\}$ ve $\{3,4,5,\dots\}$ olduğunu varsayın. O zaman dizi şöyle
başlardı:
\[
\begin{array}{r@{}rrrrrrr}
  Z_1 = \{ & \mathbf{1}, & 2, & 3, & 4, & 5, & 6, & \dots\}\\
  Z_2 = \{ &  & \mathbf{2}, &  & 4, &  & 6, & \dots\}\\
  Z_3 = \{ & 1, & 2, &  &  & 5\phantom{,} &  & \}\\
  Z_4 = \{ &  &  & 3, & \mathbf{4}, & 5, & 6, & \dots\}\\
  \vdots & & & & & \ddots
\end{array}
\]
O hâlde $\overline{Z}$, köşegen boyunca aşağı inip köşegende görünen
sayıları dışarıda bırakarak ve dizinin $j$'inci satır/sütununda boşluk olan
$j$'leri içeri alarak elde edilen kümedir. Yukarıdaki durumda $1$ ve $2$'yi
dışarıda bırakır, $3$'ü içeri alır, $4$'ü dışarıda bırakır vb.'yiz.
\end{explain}

\begin{prob}
$\Pow{\Nat}$'ın !!{nonenumerable} olduğunu bir köşegen argümanıyla gösterin.
\end{prob}

\begin{prob}\label{sfr:siz:nen:prob:f-posint}
$f\colon\PosInt\to\PosInt$ fonksiyonları kümesinin !!{nonenumerable}
olduğunu açık bir köşegen argümanıyla gösterin. Yani $f_1$, $f_2$, \dots{}
bir fonksiyonlar listesi ve her $f_i\colon\PosInt\to\PosInt$ olduğunda bu
listede bulunmayan bir $\overline{f}\colon\PosInt\to\PosInt$ fonksiyonunun
var olduğunu gösterin.
\end{prob}

\end{document}
